\documentclass[11pt]{article}
\usepackage[margin=1in]{geometry}
\usepackage{amsmath,amssymb,amsthm}
\usepackage[utf8]{inputenc}
\usepackage{listings}
\usepackage{hyperref}
\usepackage{xcolor}
\usepackage{textcomp}

\definecolor{leanblue}{rgb}{0.2,0.4,0.7}

\lstdefinelanguage{Lean}{
  keywords={theorem,lemma,def,axiom,structure,class,instance,where,by,sorry,simp,exact,intro,apply,have,let,fun,match,if,then,else,import,open,namespace,end,section,variable,inductive,noncomputable,abbrev,deriving},
  keywordstyle=\color{leanblue}\bfseries,
  comment=[l]{--},
  morecomment=[s]{/-}{-/},
  string=[b]",
  basicstyle=\ttfamily\small,
  breaklines=true,
  extendedchars=true,
  inputencoding=utf8,
  literate=
    {≥}{{$\geq$}}1
    {≤}{{$\leq$}}1
    {→}{{$\to$}}1
    {←}{{$\leftarrow$}}1
    {∧}{{$\wedge$}}1
    {∨}{{$\vee$}}1
    {¬}{{$\neg$}}1
    {∀}{{$\forall$}}1
    {∃}{{$\exists$}}1
    {⊆}{{$\subseteq$}}1
    {⊂}{{$\subset$}}1
    {∈}{{$\in$}}1
    {∉}{{$\notin$}}1
    {×}{{$\times$}}1
    {⊗}{{$\otimes$}}1
    {▸}{{$\blacktriangleright$}}1
    {⟨}{{$\langle$}}1
    {⟩}{{$\rangle$}}1
    {λ}{{$\lambda$}}1
    {α}{{$\alpha$}}1
    {β}{{$\beta$}}1
    {σ}{{$\sigma$}}1
    {θ}{{$\theta$}}1
    {ε}{{$\varepsilon$}}1
    {μ}{{$\mu$}}1
    {π}{{$\pi$}}1
    {Σ}{{$\Sigma$}}1
    {Π}{{$\Pi$}}1
    {▶}{{$\triangleright$}}1
    {⊔}{{$\sqcup$}}1
    {⊓}{{$\sqcap$}}1
    {⊥}{{$\bot$}}1
    {⊤}{{$\top$}}1
    {∅}{{$\emptyset$}}1
    {∪}{{$\cup$}}1
    {∩}{{$\cap$}}1
    {≠}{{$\neq$}}1
    {ℕ}{{$\mathbb{N}$}}1
    {₀}{{$_0$}}1
    {₁}{{$_1$}}1
    {₂}{{$_2$}}1
    {|>}{{$|{\!>}$}}2
    {·}{{$\cdot$}}1
    {—}{{---}}1
    {∘}{{$\circ$}}1
    {√}{{$\sqrt{}$}}1
    {∝}{{$\propto$}}1
    {≈}{{$\approx$}}1
    {═}{{=}}1
    {⟹}{{$\Longrightarrow$}}1,
}

\newtheorem{theorem}{Theorem}
\newtheorem{lemma}[theorem]{Lemma}
\newtheorem{definition}[theorem]{Definition}
\newtheorem{axiom_stmt}[theorem]{Axiom}

\title{Formal Verification: Hybrid Signature Scheme}
\author{Zach Kelling \and Woo Bin}
\date{December 2025}

\begin{document}
\maketitle

\begin{abstract}
We formalize the Lux hybrid signature scheme combining Ed25519 (classical), ML-DSA-65 (lattice-based), and SLH-DSA-128s (hash-based). Full verification requires all three to pass. An attacker must break all three cryptographic assumptions to forge a signature.
\end{abstract}

\section{Introduction}
The hybrid scheme provides defense in depth for Lux validator identities. In normal operation, fast mode (Ed25519 + ML-DSA) balances performance and post-quantum security. Full mode adds SLH-DSA as a hash-based backup. Classical mode is a last-resort fallback.

\section{Definitions}
\begin{definition}[HybridPK]
Hybrid public key: Ed25519 + ML-DSA + SLH-DSA components.
\end{definition}

\begin{definition}[HybridSig]
Hybrid signature: one signature from each scheme.
\end{definition}

\begin{definition}[VerifyMode]
Three modes: fast (Ed25519 + ML-DSA), full (all three), classical (Ed25519 only).
\end{definition}


\section{Theorems and Axioms}
\begin{theorem}[\texttt{full\_requires\_all\_three}]
Full verification requires all three schemes to verify: Ed25519 AND ML-DSA AND SLH-DSA.
\end{theorem}
\begin{proof}
Decompose the nested Boolean conjunction.
\end{proof}

\begin{theorem}[\texttt{fast\_requires\_both}]
Fast mode requires both Ed25519 and ML-DSA.
\end{theorem}

\begin{theorem}[\texttt{full\_implies\_fast}]
Full verification implies fast verification.
\end{theorem}
\begin{proof}
Extract the first two components from \texttt{full\_requires\_all\_three}.
\end{proof}

\begin{theorem}[\texttt{full\_implies\_classical}]
Full verification implies classical (Ed25519) verification.
\end{theorem}


\section{Lean Source}
\lstinputlisting[language=Lean]{../../formal/lean/Crypto/Hybrid.lean}

\section{References}
\noindent Source code: \url{https://github.com/luxfi/proofs/blob/main/lean/Crypto/Hybrid.lean}\\
\noindent White papers: \url{https://papers.lux.network}

\noindent Related white papers:
\begin{itemize}
  \item \texttt{lux-ethfalcon-post-quantum.tex}
\end{itemize}

\end{document}
